% Generated draft: requires mathematical and source review.
\documentclass[11pt]{article}
\usepackage[margin=25mm]{geometry}
\usepackage{fontspec}
\setmainfont{Latin Modern Roman}
\usepackage{amsmath,amssymb,mathtools,unicode-math}
\setmathfont{Latin Modern Math}
\usepackage{enumitem,needspace,xurl,hyperref}
\setlength{\parindent}{0pt}
\setlength{\parskip}{5pt}
\setlength{\emergencystretch}{3em}
\newcommand{\N}{\mathbb{N}}
\newcommand{\Z}{\mathbb{Z}}
\newcommand{\Q}{\mathbb{Q}}
\newcommand{\R}{\mathbb{R}}
\newcommand{\C}{\mathbb{C}}
\newcommand{\F}{\mathbb{F}}
\newcommand{\A}{\mathbb{A}}
\newcommand{\PP}{\mathbb{P}}
\newcommand{\E}{\mathbb{E}}
\newcommand{\eps}{\varepsilon}
\newcommand{\dd}{\,\mathrm{d}}
\begin{document}
\section*{UnsolvedMath ID: 30001714}
\textbf{Source problem code:} catalogue entry
\subsection{Definitions and mathematical statement}
The supplied record proposes a relation between a character variety and a Kac polynomial, but does not specify the data needed to define either side. Retaining its notation, the proposed statement is
\[
P_c(M_\mu;t)=t^{2d_\mu}A_\Gamma(\mu,t^2),
\qquad\text{with the asserted consequence}\qquad
A_\Gamma(\mu,q)\in\mathbb Z_{\geq 0}[q].
\]
Here \(\mathbb Z_{\geq 0}[q]\) denotes polynomials in the indeterminate \(q\) with nonnegative integer coefficients, expressing the intended coefficient-positivity assertion written as \(\mathbb N[q]\) in the record.

For clarity about the missing definitions, under the usual counting interpretation a Kac polynomial \(A_\Gamma(\alpha,q)\) is associated with a finite quiver \(\Gamma\) and a dimension vector \(\alpha\in\mathbb Z_{\geq 0}^{\Gamma_0}\), where \(\Gamma_0\) is the vertex set. At a prime power \(q\), its value counts isomorphism classes of absolutely indecomposable representations of \(\Gamma\) over \(\mathbb F_q\) of dimension vector \(\alpha\). Such a representation assigns a vector space to each vertex and a linear map to each arrow; absolute indecomposability means that it remains indecomposable after extension to an algebraic closure. The supplied record does not define \(\Gamma\), the parameter \(\mu\), or any dimension vector determined by \(\mu\), so this interpretation cannot yet be instantiated.

One common compactly supported Poincar\'e-polynomial convention for a complex algebraic variety \(X\) is
\[
P_c(X;t)=\sum_{i\geq 0}\dim_{\mathbb Q}H_c^i(X;\mathbb Q)\,t^i.
\]
This convention is stated only to expose the unresolved normalization: the record does not define \(P_c\), its grading, or whether a weight-restricted cohomological invariant is intended. It also does not define \(M_\mu\), \(d_\mu\), or their relationship. In particular, no ambient group, underlying surface, prescribed conjugacy classes, genericity conditions, or admissible range of parameters is supplied. These missing hypotheses prevent a fully specified quantified mathematical statement. The displayed formula is therefore preserved as the extracted conjectural target, pending recovery of its definitions and normalization from the cited sources.
\subsection{Short English statement}
The extracted conjecture relates compactly supported cohomology of an associated character variety to a Kac polynomial and asserts nonnegativity of its coefficients. The character variety, quiver, parameters, exponent, and cohomological normalization remain unspecified, so the formula requires source clarification before mathematical review.
\subsection{Sources}
\begin{itemize}
\item Supplied local problem registry: permanent UnsolvedMath ID \(30001714\), title \emph{Kac Polynomials and Character-Variety Cohomology}, source problem code \texttt{catalogue entry}, in the Oberwolfach Reports Open Problems collection. The record assigns proposed year 2011 and supplies no individual proposer.
\item Supplied report citation: \emph{Representation Theory of Quivers and Finite Dimensional Algebras} (2011), Oberwolfach Reports, pp.~523--608, DOI 10.4171/owr/2011/10. \url{https://doi.org/10.4171/owr/2011/10}. The supplied background describes Oberwolfach Reports as documenting mathematical workshops and their open-problem discussions.
\item Supplied literature source: Hausel--Letellier--Rodriguez-Villegas. \url{https://arxiv.org/abs/0810.2076}.
\item Supplied literature source: Mellit. \url{https://annals.math.princeton.edu/2020/192-1/p03}.
\item Supplied literature source: Oberwolfach Report 10/2011. \url{https://doi.org/10.4171/owr/2011/10}.
\item Supplied literature assessment, dated exactly \texttt{2026-08-21}: the registry labels the problem partially solved and states that the Kac-polynomial and character-variety formulas have substantial proved cases, that Mellit proved the corresponding Poincar\'e-polynomial formula, and that the full mixed-Hodge conjectural package in the extracted statement is not completely settled. This is supplied dated curation derived from the OWR material, not an independently verified assessment. The supplied background expressly requests independent verification of time-sensitive status and cited claims.
\end{itemize}
\end{document}
